\documentclass[12pt]{book}

\usepackage{graphicx}
\usepackage{eso-pic}
\usepackage{tcolorbox}
\usepackage{polyglossia}
\setdefaultlanguage[indentfirst=false,forceheadingpunctuation=false]{russian}
\setotherlanguages{english}

\RequirePackage[a4paper, headsep = 0.5 \headsep, left=2.5cm, right=2.1cm, top=2cm, bottom=2.1cm]{geometry}

\setmainfont{URW Bookman}
\newfontfamily\cyrillicfont{URW Bookman}[Script=Cyrillic, Ligatures=TeX]
\setlength{\parindent}{20pt}

\usepackage{enumitem}
\setlist[itemize, 1]{label =\raisebox{-0.3\height}{\scalebox{1.8}{\textbullet}}}

\title{Принимая отчужденность}
\author{Тодд МакГоуэн}
\date{2019}
\begin{document}

\AddToShipoutPictureBG*{\includegraphics[width=\paperwidth]{title.png}}

\null
    \thispagestyle{empty}%
    \addtocounter{page}{-1}%
    \newpage

\tableofcontents

\addcontentsline{toc}{chapter}{Введение}
\chapter*{Введение}

''Отчужденность'' звучит непривлекательно. Будто что-то разрушительное, чего лучше избегать. Никто не говорит друзьям ''Давайте отчуждаться'', наверняка это не так приятно, как основание сообщества, романтические отношения, общение с приятелями, пицца, игра в теннис или выпивка. Напротив --- мы хотим помочь подростку избежать отчужденности, а отчужденному постороннему --- почувствовать себя уютней. Левые, согласно Карлу Марксу, понимают политическую деятельность как борьбу с отчуждением. Почти всем отчужденность видится проблемой, которую надо решить.

Приговор отчужденности поспешен и забывает о ее освободительном качестве. Исторически, у отчужденности плохая репутация, но ведь именно отчужденность дает возможность выхода за пределы конкретной ситуации. Следовало бы думать о ней гораздо лучше. Увидеть ее достоинства и постараться усилить ее опыт.

Никто открыто не отстаивал отчужденность как проект, но некоторые принимают ее основой свободы. Немецкий идеализм, кульминацией которого стала философия Гегеля, говорит об отчужденности как условии самосознания. Психоанализ воспринимает отчужденность субъекта как непреклонный факт нашего существования и стремится примирить с ним. Экзистенциализм Сартра и Симоны де Бовуар тоже видит в отчужденности субъекта основу свободы. Франц Фанон объединяет все три течения --- Гегеля, психоанализ и экзистенциализм --- и формулирует антиколониальную теорию, выступая против попыток избегать фундаментальной нам всем отчужденности. Впрочем, эти мыслители никогда не используют отчужденность как основное понятие или лозунг, хотя оно объединяет эти разрозненные попытки говорить о свободе. В этой книге отчужденность выдвигается на первый план.

Отчужденность парадоксальна. Она не значит отклонения от исходной идентичности, от того, что мы когда-то имели но потеряли. Отчужденность неизменно присуща нам. Она не отдаленность от первичного, потому что первична. Мы отчуждены в своей субъектности до самой смерти. Смерть --- единственное лекарство от этой болезни\footnote{В момент смерти наша отчужденность проходит --- мы становимся совершенно чуждыми. Смертью субъект совершенно теряется в инаковости, уже не сохраняя ни малейшей частицы себя.}, так что, наверное, не такая уж это и болезнь. Отчужденность --- не то, о чем можно сожалеть или чего лучше избегать, ведь благодаря ней возможно то, что делает существование достойным продолжения --- свобода, равенство, солидарность, способность выходить за пределы ситуации.

Все творение зависит от отчужденности, она выкорчевывает нас из определенности, освобождает от конкретной ситуации. Каждый рождается в специфичном социуме, в котором на него действуют мощные силы, но не сводится к ним. Первичное отчуждение в языке дистанцирует как от других, так и внутренне: к себе теперь можно относиться как к другому. Мы отчуждены от себя и от других, и что бы мы ни делали, нам никогда не преодолеть эту отчужденность. Она основа нашего существования.

Отчужденность --- это отсутствие самоидентичности. Она отдаляет субъекта от условий, в которых он проявляется. Если я не тождествен себе, если я в противоречии с собой, то внешние силы не могут меня полностью определить. Существо, дистанцированное от самого себя, никогда не сможет быть тем, во что его толкают внешние силы. Этот внутренний раскол значит сопротивление детерминированию биологией или обществом. Именно дистанция от самой себя позволяет субъектности действовать против определяющего. Социальное и натуралистическое утрачивают свою власть на субъектом, потому что его отчужденность делает их незначительными.

Я могу не оставаться до конца своих дней жителем маленького американского городка на Среднем Западе, могу бросить вызов определяющим факторам этой ситуации, и все потому что никогда не отождествлен с положением в котором родился или нахожусь. Это дает мне возможность отнестись к ситуации критически, вплоть до отказа от всех идей и ценностей, которые возложило на меня общество. Я могу быть воспитан с энтузиазмом к капиталистическому строю, но стать марксистом. Если меня заставляют быть женственным, я могу стать мужественным. Полюбив Средний Запад Америки, я, возможно, стану работать в Аргентине. Моя неспособность быть именно тем, кто я есть --- моя не-однозначность --- вот что позволяет мне выйти из исходной ситуации. И даже если я остаюсь в ней и принимаю ее, я все равно отношусь к ней как отчужденный. Можно быть чужим и у себя дома. У меня, как и у всех, есть внутренняя чужеродность, мешающая прямому усвоению диктата общества.

Без этой внутренней дистанции человек не был бы субъектом, не смог бы соотноситься с собой. Субъектность порождена и определена этим расколом. Отделенность от самих себя делает нас нами --- неотделимыми от чуждого в нас\footnote{Славой Жижек в \textit{Добавочном Наслаждении} подчеркивает, что отчуждение конститутивно. Мы не целостны изначально, чтобы \textit{затем} отчуждаться. Мы возвращаемся к изначальному через отчуждение. Поэтому ''вне отчужденности никакой жизни нет, нет никакой позитивной основы отчужденности'' (Slavoj Žižek, \textit{Surplus Enjoyment: A Guide for the Non-Perplexed} (London: Bloomsbury, 2022), 39). Верить, что мы были до и вне отчуждения --- ностальгировать по истоку, которого никогда не существовало.}. Чуждое дает нам нашу уникальность.

Эта внутренняя чуждость и есть то, что Зигмунд Фрейд назовет бессознательным\footnote{Для описания своей концепции Фрейд умышленно выберет \textit{unbewußt}, а не распространенный в то время \textit{bewußtlos}. \textit{Bewußtlos} значил бы ''без сознания'', \textit{unbewußt} же значит противоположное сознанию, чуждое его структуре.}. Бессознательное ставит барьер как между субъектом и ситуацией, так и внутри субъекта. Только отчужденные обладают бессознательным. Вернее, чуждый статус бессознательного значит, что это оно владеет нами. Это иное внутри диктует субъекту его действия, направляет усилия и создает убеждения. Бессознательное ориентирует субъектность, потому что выражает ее самораскол\footnote{В психике Фрейд даст бессознательному структурный приоритет. Ускользая от контроля, оно гораздо больше влияет на нас, чем сознательное мышление: ''Бессознательное --- вот истинная психическая реальность'' (Sigmund Freud, \textit{The Interpretation of Dreams}, trans. James Strachey, in The Complete Psychological Works of Sigmund Freud, vol. 5, ed. James Strachey (London: Hogarth Press, 1953 [1900]), 651). Фрейд написал это на заре своей теории, впоследствии он откажется от некоторых своих идей, но никогда не усомнится в приоритете бессознательного.}.

Отчуждение субъектов от себя видно, когда бессознательное заставляет их идти наперекор своим же интересам. Кто-то объедается пончиками и кексами, вместо того, чтобы есть яблоки и брокколи и предотвратить ишемическую болезнь. Другой вместо комфортного, неконфликтного партнера выбирает сложного, болезненного человека, с которым постоянно ссорится. Продолжает искать и употреблять все новые дурманящие вещества, хотя прекрасно понимает, как они опасны и вредны. Хочет профессионального роста, но раз за разом выступает против начальника. Повсюду субъекта поджидают саморазрушительные, ускользающие от сознательного контроля действия.

Каждый сосуществует с идущей вразрез его желаниям силой. Каждый видит, как сам мешает достижению своих целей, как не может контролировать себя. Отсутствие контроля изнурительно, но также основа способности субъекта становиться иным. Верховенство чуждого --- бессознательного --- открывает возможность свободы. Свобода напрямую зависит от подрыва сознательных желаний бессознательными импульсами.

О бессознательном обычно думают как о принуждающем, забирающем свободу. Мы не обдумываем бессознательные действия и не чувствуем, что они выбраны нами. Но свобода --- не плод сознательных размышлений. Это прерывание решимости, определенности. И именно это делает бессознательное. Оно перебивает социальные и психологические определяющих факторы, открывает разрыв в необусловленную ими альтернативную возможность. Избегая сознательного контроля, бессознательное становится точкой свободы субъекта. Мы свободны благодаря своей раздвоенности.

Самотождественности лишена не только субъективность. Ничто не самотождественно, иначе не смогло бы меняться во времени и перемещаться в пространстве\footnote{Жан-Поль Сартр утверждает, что только субъекты отчуждены, а остальные существа тождественны себе. Он радикально различает субъекта, который для-себя, и остальное бытие, которое в-себе: ''Бытие есть. Бытие --- в-себе. Бытие есть то, что оно есть.'' (Jean-Paul Sartre, \textit{Being and Nothingness}, trans. Hazel E. Barnes (New York: Washington Square Press, 1956 [1943]), 29). Сартр здесь не задумывается, как субъект мог стать отчужденным. Будь бытие лишь тем, что оно есть, субъект не возник бы, не проявился из такого чистого царства самотождественности.}. Это касается и животных и растений, и гор и рек. Будь сущность тем, что она есть, она бы никогда не изменилась. Субъект же превращает свою нетождественность в определяющую его структуру, такой процесс и есть отчуждение.

Первое свидетельство отчуждения --- имя. Ни один ребенок не дает его себе сам\footnote{Мама говорит, что я сам выбрал имя: она положила в шляпу несколько листочков бумаги с именами, и я схватил листочек с именем ''Тодд''. Такой метод не избегает навязывания, ведь все равно контролируется взрослым: он выбирает возможные имена, вырезает и кладет в шляпу листочки с ними. Заключенное в присвоении имени отчуждение неизбежно. Затем мама призналась, что ''Тодд'' было написано на каждом из листочков. Впрочем, признание могло быть следствием тогда уже развившейся у нее болезни Альцгеймера, но так или иначе, оно показывает определяющую роль внешнего в присвоении имени.}. Происходит акт насилия, родители или другие взрослые навязывают ребенку произвольное имя. И даже если он потом воспринимает его своим, оно остается чуждым, сошедшим на него извне. Родители могут называть мальчика Эрнестом и надеется, что он вырастет честным и прямодушным. Могут называть девочку Грейс и думать, что она будет приличной и изящной. Иногда используют имена художников, таких как Лэнгстон или Дэшил, пытаясь ассоциацией приобщить детей к их гениальности. Дают имя, связанное с семейной историей, чтобы оно напоминало ребенку о его происхождении, определяло его поведение. Но и ничего не значащее имя остается чуждым названному означающим. Он не выбирал его, и воспринимает его своим только через отчуждение.

Имя это наша дистанция до себя. В каждом языке есть способ идентифицировать себя именем: я могу сказать ''Меня зовут Тодд'', или ''Ich heiße Todd'', или ''Me llamo Todd''. Всякое такое высказывание провозглашает тождество, и этим выявляет внутреннюю дистанцию между говорящим и именем. Так оно выдает, что я и Тодд --- не одно и то же\footnote{В логике первого порядка утверждение тождества исключает различие, потому что различие значило бы категориальную ошибку: смешение идентификации и предикации. На этом замечании Бертран Рассел строит свою критику Гегеля. Гегель ошибается --- утверждает он --- когда говорит об идентичности в различии, об этой фундаментальной отчужденности субъекта, ведь, например, ''Скотт --- автор ''Уэверли'''' предполагает скорее идентификацию, чем предикацию. Критику этой позиции Рассела и защиту Гегеля см. в \textit{Emancipation After Hegel: Achieving a Contradictory Revolution} (New York: Columbia University Press, 2019).}. Самотождественное существо не смогло бы так сказать.

Даже если я переименую себя, имя останется чуждым. Казалось бы, я называю себя сам, не полагаюсь на данное другими. Разве этого недостаточно для преодоления отчужденности? Нет. Называя себя иначе, я действую исходя их доступных возможностей: выбираю из используемых другими имен или комбинирую их. Выбор имени всегда предполагает долг перед внешним. Мое новое имя не только принадлежит другим, но и выбрано с учетом их впечатления. Откуда его не возьми, оно несет печать самораскола.

Неизбежность отчужденности дает некоторую уверенность. И вместо того, чтобы страдать от саморазделения до самой смерти, мы можем манипулировать им и даже усиливать его. Оставление тщетных попыток преодолеть отчужденность, ее принятие, трансформирует ее. Из бремени она становится освобождением от определенности. Переосмысление отчужденности --- важнейшая политическая задача, особенно в эпоху вездесущих схем его якобы преодоления.

Об отчужденности плохо думают, потому что воображают скрывающуюся за ней завершенность. Где-то в затерянном прошлом, в Эдеме, или в утопическом будущем, нас ждет самотождественность, полнота, совершенство. Она всегда ускользает и тем соблазняет еще больше. Такое лишенное всякого отчуждения состояние очень привлекает нас, но лишь пока остается воображаемым. Малейшая попытка его очертить, конкретизировать, воплотить, выдает его полную несостоятельность\footnote{Никто не считает \textit{Возвращенный рай} лучшей поэмой, чем \textit{Потерянный рай}. Отчужденность приятна, потому что ее неприятно пытаться преодолевать. Мильтон невольно демонстирует это разницей в качестве поэм и --- более того --- изображает неотчужденных Бога, Христа, ангелов в \textit{Потерянном раю} совсем не такими привлекательными, как отчужденного Сатану.}.

Нам следует мыслить отчужденность вне ассоциаций с полнотой или самотождественностью. Мы отчуждены, но не от чего-то могущего дополнить нас. Мы и не теряли ничего такого. Помыслив нехватку вне связи с целостностью, мы поменяем политическую значимость отчужденности. Отчужденность не определяется через не-отчужденность.

Отчужденность --- источник всех человеческих страданий, но оно и источник всякого удовлетворения. Отчужденность предоставляет пространство действию вопреки, без него субъект неотличим от сил, породивших его. Заниматься сексом, есть торт или обсуждать немецкий идеализм --- ничего из этого не возможно без отчужденности, потому что без нее нет субъекта.

Если что и отличает говорящих существ от прочих, так это их осознание отсутствия самоидентичности. Через язык отчужденность осознается субъектами неизменно присущей им. Язык создает дистанцию, и это меняет отношение существа к отчужденности. Субъекты означающего могут осознать и принять себя отчужденными, могут найти удовлетворение в этом отсутствии самотождественности, и только отказ рассматривать отчужденность конститутивной составляющей мешает этому.

Чаще мы думаем о своих экзистенциальных и политических проектах как о преодолевающих отчужденность. За редкими исключениями, мы стремимся дать отчужденности бой, искоренить ее. Психотерапевты воображают, что помогают пациентам обрести самоидентификацию. Революционеры стремятся построить общество, в котором жили бы в гармонии с другими. Экоактивисты работают на формой существования, которая была бы в гармонии с природой\footnote{Впрочем, не всегда. Например, эколог Мюррей Букчин, будучи гегельянцем, говорит, что для экологического проекта необходимо принять отчужденность. См. Murray Bookchin, \textit{The Philosophy of Social Ecology: Essays on Dialectical Naturalism} (Chico: AK Press, 2022 [1990]).}. Разительное разнообразие проектов не затрагивает их сути --- избавления от чуждого. Отчужденность непрестанно видится тем, с чем следует бороться, хотя сама борьба возможна только благодаря ней.

Эта книга предлагает совершенно иной взгляд на отчужденность: усилия в ее преодолении значат не радикальный ответ ситуации, а неспособность увидеть конституирующую роль отчужденности. Нам следует оставить попытки преодолеть ее. Отчужденность должна стать нашей экзистенциальной и политической программой. Отчуждение --- это освобождение.

\chapter{Беспокойный субъект}

\section{Продажная идентичность}

Субъект не то же, что присвоенная ему символическая идентичность. Невозможно понять отчуждение, не видя этого базового различия --- между субъектностью и идентичностью. Идентичность складывается из социальных позиций: работа, семейный статус, религиозные пристрастия, политические предпочтения, национальность и т.д.. Как бы легко я себя в них не видел, я никогда не отождествляюсь с ними полностью. Этот недостаток и есть субъектность. Ни одна из идентичностей не может полностью совпасть с ней, как не подгоняй.

Субъектность возникает в отчуждении, в разрыве между субъектом и принимаемой им символической идентичностью. Субъект не способен стать символической идентичностью, потому что соотнесение с ней значит дистанцию от нее. Всегда есть разрыв между мной как субъектом и мной как личностью. Само усилие, порыв соответствовать идентичности выражает расхождение с ней\footnote{Жан-Поль Сартр показывает это в своем рассказе об официанте, ведущем себя так, будто он официант. Пытаясь полностью воплотить символическую идентичность, он неизбежно терпит неудачу: ''Но посмотрите на этого официанта. Его движения быстрые и точные. Чуть слишком точные и чуть слишком быстрые. Он идет к посетителям чуть слишком быстрым шагом, наклоняется вперед чуть слишком нетерпеливо, голос и глаза его выражают чуть слишком заботливый интерес к заказу клиента'' (Jean-Paul Sartre, Being and Nothingness, trans. Hazel E. Barnes (New York: Washington Square Press, 1956 [1943]), 59). Для Сартра вера в то, что человек и есть символическая идентичность --- самообман, и даже обманщик не может преодолеть разрыв между ними. Слишком высокая точность выдает несовершенство.}. Ведь сводись субъектность к символической идентичности, самоидентификация и не понадобилась.

Субъектность неопределима. Кто мы как субъекты --- вопрос без ответа, хотя мы и пытаемся ответить на него: мужчина, женщина, китаец, итальянец, индус, мусульманин, юрист, дантист. Когда я чувствую, что теряю контроль, плотность существующего субъекта, я напоминаю себе, что я американец, или что я сын, или что я профессор --- напоминаю себе, что я есть. Но никакие усилия не снимают дистанцию между субъектностью и идентичностью\footnote{Чем больше прикладывается усилий соответствовать, тем явственней несоответствие. Усилие значит отсутствие того, на что направлено. Если бы я на самом деле соответствовал своей идентичности профессора, мне не пришлось бы стараться одеваться как профессор, говорить как профессор и ходить как профессор.}. Все попытки соответствовать проваливаются\footnote{Субъект не моряк Попай и не может сказать ''Я таков, каков я есть'' (''I Yam what I yam and that's all what I yam'').}.

Символическая идентичность и маскирует, и маркирует вопрос. Путаница неизбежна, потому что субъектности без присвоенной ей символической идентичности не бывает. Идентичность неизбежна, только она ориентирует нас в социальной ситуации.

И в то же время идентичность --- бегство от субъектности. Мы ищем передышку в борьбе за безопасность, надеемся укрыться в идентичности, но она всегда не то, что мы искали. Отношение к ней определяет наше политическое бытие.

Мы можем противостоять символической идентичности отчужденной субъектностью. Кажется, что это бесполезно, но это политический акт --- первый и самый важный, потому что позволяет действовать за пределами символического.

Погружаясь в символическую идентичность, мы избегаем этого противостояния. Избегаем политических вопросов, на которые, кажется, отвечает идентичность. Символическая идентичность значит избегание проблемы свободы. Этот момент подчеркивает Аленка Зупанчич в \textit{Что есть секс?}: ''(освободительная) политика начинается с ''утраты идентичности'', утраты, в которой нет ничего прискорбного''\footnote{Alenka Zupančič, What IS Sex? (Cambridge: MIT Press, 2017), 36}. Идентичность деполитизирует. Всякая идентичность конформна.

Какой бы трансгрессивной не казалась идентичность, субъект, когда ее принимает --- или пытается принять --- подчиняется силам установившей ее социальной власти. Бессознательное желание субъекта всегда проскакивает мимо требований социального авторитета, и в поисках одобрения субъект моделирует себя идентификацией по требуемому образцу. Бессознательное действует вне авторитета, а символическая идентичность ищет его признания. Субъект принимает символическую идентичность, чтобы ее признал социальный авторитет, поэтому она сразу уже им ограничена. И диссонанс между идентичностями не выводит за эти границы, в основе каждой из них все еще лежит социальное признание. Сколь бы радикальной не воображалась идентичность, через нее в конечном счете ищется чье-то одобрение.

Политическое различие между субъектностью и символической идентичностью разоблачается всякий раз, когда возникает скользкая, сложная социальная проблема. Вопрос про отношение к смертной казни, например, провоцирует положиться на символическую идентичность, дать заготовленный ответ, не обдумав все варианты. Принять точку зрения Католицизма и ответить как католик --- что против смертной казни. Упомянуть, что родители воспитали меня в строгом соответствии с законом Десяти Заповедей, так что смертная казнь --- выражение справедливости. Решить, что я либерал, и что, значит, смертная казнь --- слишком суровое наказание. В каждом варианте некоторая символическая идентичность помогает отставить бремя субъектности. Идентичность предлагает способы решения экзистенциальной проблемы субъектности. В любой непонятной ситуации я идентифицируюсь с определенной символической позицией, и получаю ясный ответ.

Переход от субъектности к идентичности закрывает вопросы, на которые нет готовых ответов. Идентичность ориентирует и показывает, как реагировать, поэтому мы ищем в ней убежище. Но как субъект, я должен ответить на экзистенциальные вопросы вне ее руководства. Соответствие ответа некоторой идентичности не оправдывает его.

Символическая идентичность манит в бегство от субъектности\footnote{Принятая идентичность сама может стать в тягость, если я не привык ей соответствовать.}. Поиск идентичности, которая бы полностью нам соответствовала --- это попытка вылечить отчужденность. Субъект видится сводимым к идентичности, когда она затмевает субъектность. Субъект всегда в разладе с собой, но идентичность --- даже если находится в постоянном движении --- создает мираж целостности. Мы можем оставаться в рамках вопроса о субъектности только признавая всякую идентичность неудачной. Такое признание возможно, потому что субъект никогда не соответствует закрепленной за ним идентичности. Контраст с идентичностью выводит субъекта на первый план.

Речь вносит дистанцию в говорящего. Субъектность и есть эта дистанция, эта неспособность совершенно быть собой. Отношение субъекта к самому себе опосредуется означающим, внешним. Означающее и образует форму символической идентичности, и выявляет раскол, заставляющий субъекта искать ответ на вопрос, чем он является. Всякое означающее диалектично: подчинившись ему, существо обретает способность мыслить отрицание себя. В \textit{Прочти мое желание} Джоан Копжек утверждает: ''обозначение неизбежно порождает сомнение, возможность собственного отрицания; научает думать об аннигиляции, полномасштабном разрушении всей нашей означаемой реальности''\footnote{Joan Copjec, \textit{Read My Desire: Lacan Against the Historicists} (Cambridge: MIT Press, 1994), 54–55.}. Без возможности отрицания --- возможности отринуть не только присвоенное значение, но целый означаемый универсум --- нет субъектности. Для отчужденного все данное уязвимо сомнению и может быть отвергнуто, даже утверждать свое означающее он может только через нехватку самотождественности.

Идентифицируя себя как немец, субъект показывает свою нетождественность с немецкостью. Но любая, даже самая прямая идентификация с немецкостью, только лишний раз демонстрирует ее чуждость субъекту. Если бы немецкость действительно не отличалась от него, акт идентификации был излишним и --- более того --- невозможным. Утверждение повторения свидетельствует о различии. Обычно мы не осознаем, как самоидентификация означающими самодистанцирует. Мы используем чужие слова, чтобы описать себя, никто не изобретает означающие. Акт утверждения совпадения --- признак расхождения. Это печать отчужденности.

Марин Ле Пен, лидер французской партии ''Национальное объединение'' заинтересована соответствовать тому, что значит быть француженкой. Для нее Франция --- католическая страна, и должна быть экономически независимой. Традиции Франции, включая язык, отделяют ее от остальной Европы и от мусульманских стран. Но как бы Ле Пен не заботилась о своей французской идентичности, она отличается от нее. Сама ее артикуляция этой идентичности вносит зазор. Никто не привержен традициям Франции так, как Ле Пен, но даже ей приходится трудиться, чтобы укоренить себя в своей французскости.

Самоописание раскрывает самоудаленность. Если бы субъект был идентичен французскости, никакое самоописание не потребовалось бы. Означающее идентичности указывает, что субъект не тот, кем себя называет, что должен отыгрывать свою идентичность, а не просто быть ею. Идентичность субъекта всегда вне его самого. Повторяющимися актами идентификации он принимает и поддерживает ее. Никто не получает идентичность бесплатно.

Множество научно-фантастических фильмов фокусируются на различии между субъектом и его идентичностью. В \textit{Бегущем по лезвию} (Ридли Скотт, 1981), \textit{Вспомнить все} (Пол Верховен, 1990) и \textit{Темный город} (Алекс Пройас, 1998) мы видим, как эта дистанция проявляется в повествовании\footnote{В \textit{Бегущем по лезвию} Рик Декард (Харрисон Форд) --- охотник за головами --- преследует беглых репликантов, искусственно созданных существ, которым для приручения вживили воспоминания. Но понимая, что его собственная идентичность столь же ложна, как идентичности репликантов, Декард освобождается от нее. Он, может быть, тоже репликант, но даже если не так --- в символической идентичности нет надежности. Фильм заканчивается тем, что Декард вступает в романтические отношения с женщиной, зная, что она репликант.}.

\end{document}
